\documentclass[UTF8,12pt]{ctexart}

\usepackage[a4paper,margin=2.4cm]{geometry}
\usepackage{amsmath,amssymb}
\usepackage{enumitem}
\usepackage{hyperref}
\usepackage{xcolor}
\usepackage{fvextra}
\usepackage{titlesec}
\usepackage{fancyhdr}
\setmonofont{DejaVu Sans Mono}
\setlength{\headheight}{15pt}

\hypersetup{
  colorlinks=true,
  linkcolor=blue,
  urlcolor=blue,
  citecolor=blue
}

\setlist[itemize]{itemsep=2pt,topsep=4pt}
\setlist[enumerate]{itemsep=2pt,topsep=4pt}
\DefineVerbatimEnvironment{LeanCode}{Verbatim}{breaklines=true,breakanywhere=true,fontsize=\small,frame=single,rulecolor=\color{gray}}
\DefineVerbatimEnvironment{Tree}{Verbatim}{breaklines=true,breakanywhere=true,fontsize=\small}

\pagestyle{fancy}
\fancyhf{}
\lhead{Lean 证明拆分规则说明文档}
\rhead{\thepage}

\title{Lean 证明拆分规则说明文档}
\author{}
\date{}

\begin{document}
\maketitle
\tableofcontents
\newpage

\section{文档目标}

当我们遇到一个复杂 Lean theorem 时，应该如何把它拆成更小、更容易证明的子目标？

Lean 证明看起来复杂，很多时候是因为目标里嵌套了很多逻辑结构，例如：
\begin{itemize}
  \item \(A \land B\)
  \item \(A \leftrightarrow B\)
  \item \(A \to B\)
  \item \(\forall x, P\ x\)
  \item \(\exists x, P\ x\)
  \item \(A \lor B\)
\end{itemize}

一个复杂 theorem 往往就是这些结构不断嵌套在一起。

所以 Lean 证明拆分的核心思想是：

\begin{quote}
先看当前目标和已知条件的形状，再根据形状选择对应的拆分规则，递归地把大目标变成小目标。
\end{quote}

\section{Lean 证明中最重要的两个对象}

Lean 证明过程中，我们一直面对两类东西：
\begin{enumerate}
  \item 当前目标；
  \item 已知条件。
\end{enumerate}

\subsection{当前目标}

Lean 中 \(\vdash\) 后面的内容就是当前要证明的目标。

例如：

\begin{LeanCode}
A B : Prop
hA : A
hB : B
⊢ A ∧ B
\end{LeanCode}

意思是：
\begin{itemize}
  \item 已知 \(A\)、\(B\) 是命题；
  \item 已知 \texttt{hA} 是 \(A\) 的证明；
  \item 已知 \texttt{hB} 是 \(B\) 的证明；
  \item 现在要证明 \(A \land B\)。
\end{itemize}

所以看到 Lean 状态时，首先要看：

\begin{quote}
\(\vdash\) 后面的目标长什么样？
\end{quote}

\subsection{已知条件}

例如：

\begin{LeanCode}
hA : A
hB : B
h : A → B
\end{LeanCode}

意思是：
\begin{itemize}
  \item \texttt{hA} 是 \(A\) 的证明；
  \item \texttt{hB} 是 \(B\) 的证明；
  \item \texttt{h} 是一个从 \(A\) 推出 \(B\) 的规则。
\end{itemize}

很多时候目标本身不好拆，但已知条件可以拆。所以除了看当前目标，也要看上下文里有没有可以利用的已知条件。

\section{复杂 Lean 证明的基本拆分思路}

遇到一个复杂 theorem，可以按下面的顺序思考：

\begin{enumerate}
  \item 先看目标最外层结构。
  \item 如果目标是 \(\to\) 或 \(\forall\)，先 \texttt{intro}。
  \item 如果目标是 \(\land\) 或 \(\leftrightarrow\)，用 \texttt{constructor} 拆成多个子目标。
  \item 如果目标是 \(\exists\)，先找 witness。
  \item 如果目标是 \(\lor\)，选择证明左边或右边。
  \item 如果已知条件是复合结构，把已知条件拆开。
  \item 如果有定理可以推出当前目标，用 \texttt{apply} 反推。
  \item 如果目标可以化简，用 \texttt{simp} / \texttt{rw} / \texttt{unfold}。
  \item 如果目标涉及 \texttt{Nat} / \texttt{List} / 递归结构，用 \texttt{induction} 或 \texttt{cases}。
  \item 如果目标已经很小，用 \texttt{exact} / \texttt{rfl} / \texttt{contradiction} 关闭。
\end{enumerate}

这套思路可以理解成：

\begin{Tree}
复杂目标
 ↓
先拆最外层逻辑结构
 ↓
得到若干中等目标
 ↓
继续递归拆每个中等目标
 ↓
得到很多小目标
 ↓
用已有条件、定理、化简、归纳等方法关闭小目标
\end{Tree}

\section{第一类规则：目标本身可以拆}

这一类规则看的是 \(\vdash\) 后面的目标。

\subsection{目标是 \(A \land B\)}

目标形状：

\begin{LeanCode}
⊢ A ∧ B
\end{LeanCode}

意思是：要同时证明 \(A\) 和 \(B\)。

拆分方式：
\begin{itemize}
  \item 子目标 1：\(\vdash A\)
  \item 子目标 2：\(\vdash B\)
\end{itemize}

Lean 常用写法：

\begin{LeanCode}
constructor
\end{LeanCode}

例子：

\begin{LeanCode}
theorem demo
  (A B : Prop)
  (hA : A)
  (hB : B) :
  A ∧ B := by
  constructor
  · exact hA
  · exact hB
\end{LeanCode}

任务树：

\begin{Tree}
A ∧ B
├── A
└── B
\end{Tree}

只要看到目标最外层是 \(\land\)，就可以先拆成左右两个子目标。

\subsection{目标是 \(A \leftrightarrow B\)}

目标形状：

\begin{LeanCode}
⊢ A ↔ B
\end{LeanCode}

意思是：\(A\) 和 \(B\) 等价。

证明等价命题需要证明两个方向：
\begin{enumerate}
  \item \(A \to B\)
  \item \(B \to A\)
\end{enumerate}

拆分方式：
\begin{itemize}
  \item 子目标 1：\(\vdash A \to B\)
  \item 子目标 2：\(\vdash B \to A\)
\end{itemize}

Lean 常用写法：

\begin{LeanCode}
constructor
\end{LeanCode}

例子：

\begin{LeanCode}
theorem demo
  (A B : Prop)
  (h1 : A → B)
  (h2 : B → A) :
  A ↔ B := by
  constructor
  · exact h1
  · exact h2
\end{LeanCode}

任务树：

\begin{Tree}
A ↔ B
├── A → B
└── B → A
\end{Tree}

所以遇到 \(\leftrightarrow\)，不要想着一次证明“等价”，而是先拆成两个方向。

\subsection{目标是 \(A \to B\)}

目标形状：

\begin{LeanCode}
⊢ A → B
\end{LeanCode}

意思是：如果 \(A\) 成立，那么 \(B\) 成立。

注意，\(A \to B\) 不是拆成：
\begin{itemize}
  \item 证明 \(A\)；
  \item 证明 \(B\)。
\end{itemize}

而是：先假设 \(A\) 成立，然后在这个假设下证明 \(B\)。

拆分方式：

\begin{Tree}
原目标：⊢ A → B
变成：
hA : A
⊢ B
\end{Tree}

Lean 常用写法：

\begin{LeanCode}
intro hA
\end{LeanCode}

例子：

\begin{LeanCode}
theorem demo
  (A B : Prop)
  (h : A → B) :
  A → B := by
  intro hA
  exact h hA
\end{LeanCode}

任务树：

\begin{Tree}
A → B
└── 在 hA : A 的条件下证明 B
\end{Tree}

所以 \(\to\) 是一种单分支推进，不是并行拆分。\texttt{intro hA} 不是证明了 \(A\)，而是临时假设 \(A\) 成立，然后要求你在这个假设下证明 \(B\)。

\subsection{目标是 \(\forall x, P\ x\)}

目标形状：

\begin{LeanCode}
⊢ ∀ x, P x
\end{LeanCode}

意思是：对所有 \(x\)，都要证明 \(P\ x\)。

证明方式是：任取一个 \(x\)，然后证明 \(P\ x\)。

拆分方式：

\begin{Tree}
原目标：⊢ ∀ x, P x
变成：
x : α
⊢ P x
\end{Tree}

Lean 常用写法：

\begin{LeanCode}
intro x
\end{LeanCode}

例子：

\begin{LeanCode}
theorem demo
  (P : Nat → Prop)
  (h : ∀ n, P n) :
  ∀ n, P n := by
  intro n
  exact h n
\end{LeanCode}

任务树：

\begin{Tree}
∀ x, P x
└── 任取 x，证明 P x
\end{Tree}

\(\forall\) 和 \(\to\) 很像，都是先 \texttt{intro}，把目标向里面推进一层。

\subsection{目标是 \(\exists x, P\ x\)}

目标形状：

\begin{LeanCode}
⊢ ∃ x, P x
\end{LeanCode}

意思是：存在一个 \(x\)，使得 \(P\ x\) 成立。

证明存在性需要两步：
\begin{enumerate}
  \item 给出一个具体 witness \(w\)；
  \item 证明 \(P\ w\)。
\end{enumerate}

Lean 常用写法：

\begin{LeanCode}
use w
\end{LeanCode}

例子：

\begin{LeanCode}
theorem demo : ∃ n : Nat, n + 1 = 4 := by
  use 3
  norm_num
\end{LeanCode}

这里 \texttt{use 3} 的意思是：我选择 \(3\) 作为那个存在的 \(n\)。然后目标会变成：

\begin{LeanCode}
⊢ 3 + 1 = 4
\end{LeanCode}

这个目标可以用 \texttt{norm\_num} 解决。

任务关系：

\begin{Tree}
∃ x, P x
├── 找 witness w
└── 证明 P w
\end{Tree}

注意，目标是 \(\exists x, P\ x\) 时，我们是在主动构造一个存在性证明，所以必须给出具体 witness。

\subsection{目标是 \(A \lor B\)}

目标形状：

\begin{LeanCode}
⊢ A ∨ B
\end{LeanCode}

意思是：要证明 \(A\) 或 \(B\) 至少一个成立。

证明 \(A \lor B\) 不需要两边都证明，只要证明其中一边。

如果证明左边：

\begin{LeanCode}
left
\end{LeanCode}

如果证明右边：

\begin{LeanCode}
right
\end{LeanCode}

例子：

\begin{LeanCode}
theorem demo
  (A B : Prop)
  (hA : A) :
  A ∨ B := by
  left
  exact hA
\end{LeanCode}

任务关系：

\begin{Tree}
A ∨ B
├── 路线 1：证明 A
└── 路线 2：证明 B
\end{Tree}

只需要一条路线成功。

\section{第二类规则：已知条件可以拆}

很多时候目标本身不好拆，但上下文里的已知条件可以拆。这类规则看的是 \(\vdash\) 前面的已知条件。

\subsection{已知 \(h : A \land B\)}

已知条件形状：

\begin{LeanCode}
h : A ∧ B
\end{LeanCode}

意思是：已经知道 \(A\) 和 \(B\) 都成立。

可以拆成：

\begin{LeanCode}
hA : A
hB : B
\end{LeanCode}

Lean 常用写法：

\begin{LeanCode}
rcases h with ⟨hA, hB⟩
\end{LeanCode}

例子：

\begin{LeanCode}
theorem demo
  (A B : Prop)
  (h : A ∧ B) :
  A := by
  rcases h with ⟨hA, hB⟩
  exact hA
\end{LeanCode}

直觉：如果手里有 \(A \land B\)，就可以分别拿到 \(A\) 和 \(B\)。

\subsection{已知 \(h : \exists x, P\ x\)}

已知条件形状：

\begin{LeanCode}
h : ∃ x, P x
\end{LeanCode}

意思是：已经知道存在某个 \(x\) 满足 \(P\ x\)。

这个已知条件可以理解成一个“包裹”，里面装了两样东西：
\begin{enumerate}
  \item 某个具体对象 \(w\)；
  \item 证明这个对象满足 \(P\) 的证据 \(hw : P\ w\)。
\end{enumerate}

所以可以拆出：

\begin{LeanCode}
w : α
hw : P w
\end{LeanCode}

Lean 常用写法：

\begin{LeanCode}
rcases h with ⟨w, hw⟩
\end{LeanCode}

例子：

\begin{LeanCode}
theorem demo
  (P : Nat → Prop)
  (h : ∃ n : Nat, P n) :
  ∃ n : Nat, P n := by
  rcases h with ⟨w, hw⟩
  use w
  exact hw
\end{LeanCode}

解释：

\begin{LeanCode}
h : ∃ n : Nat, P n
\end{LeanCode}

表示：已经知道存在某个自然数 \(n\)，使得 \(P\ n\) 成立。

执行：

\begin{LeanCode}
rcases h with ⟨w, hw⟩
\end{LeanCode}

以后，Lean 得到：

\begin{LeanCode}
w : Nat
hw : P w
\end{LeanCode}

也就是：
\begin{itemize}
  \item \(w\) 是那个被证明存在的自然数；
  \item \(hw\) 是 \(w\) 满足 \(P\) 的证明。
\end{itemize}

如果当前目标还是：

\begin{LeanCode}
⊢ ∃ n : Nat, P n
\end{LeanCode}

那么可以：

\begin{LeanCode}
use w
\end{LeanCode}

意思是：我就用刚刚从 \(h\) 里面拆出来的 \(w\)，作为这个存在性目标的 witness。然后目标变成：

\begin{LeanCode}
⊢ P w
\end{LeanCode}

而我们已经有：

\begin{LeanCode}
hw : P w
\end{LeanCode}

所以：

\begin{LeanCode}
exact hw
\end{LeanCode}

直接完成。

最重要的区别是：
\begin{itemize}
  \item 目标是 \(\exists x, P\ x\)：你要主动给出一个 witness。
  \item 已知是 \(h : \exists x, P\ x\)：witness 已经藏在 \(h\) 里面，你要把它取出来。
\end{itemize}

\subsection{已知 \(h : A \lor B\)}

已知条件形状：

\begin{LeanCode}
h : A ∨ B
\end{LeanCode}

意思是：已知 \(A\) 或 \(B\) 至少一个成立。

但你不知道到底是哪边成立，所以必须分情况讨论：
\begin{itemize}
  \item 情况 1：假设 \(A\) 成立，证明目标；
  \item 情况 2：假设 \(B\) 成立，证明目标。
\end{itemize}

Lean 常用写法：

\begin{LeanCode}
cases h with
| inl hA =>
  ...
| inr hB =>
  ...
\end{LeanCode}

例子：

\begin{LeanCode}
theorem demo
  (A B C : Prop)
  (h : A ∨ B)
  (hAC : A → C)
  (hBC : B → C) :
  C := by
  cases h with
  | inl hA =>
    exact hAC hA
  | inr hB =>
    exact hBC hB
\end{LeanCode}

任务树：

\begin{Tree}
h : A ∨ B, goal C
├── 情况 1：hA : A ⊢ C
└── 情况 2：hB : B ⊢ C
\end{Tree}

注意：
\begin{itemize}
  \item 目标是 \(A \lor B\)：只需要证明一边。
  \item 已知是 \(h : A \lor B\)：必须两种情况都处理。
\end{itemize}

\section{第三类规则：用已有定理反推目标}

有时目标本身看不出怎么拆，但是已知条件里有一个定理可以推出目标。这时要用 \texttt{apply}。

\subsection{已知 \(h : A \to B\)，目标是 \(B\)}

当前状态：

\begin{LeanCode}
h : A → B
⊢ B
\end{LeanCode}

可以用 \(h\) 反推目标：要证明 \(B\)，只需要先证明 \(A\)。

Lean 常用写法：

\begin{LeanCode}
apply h
\end{LeanCode}

目标变化：

\begin{Tree}
原目标：⊢ B
新目标：⊢ A
\end{Tree}

例子：

\begin{LeanCode}
theorem demo
  (A B : Prop)
  (h : A → B)
  (hA : A) :
  B := by
  apply h
  exact hA
\end{LeanCode}

直觉：我有 \(A \to B\)。所以为了证明 \(B\)，可以转而证明 \(A\)。

\subsection{多前提 apply}

如果已知：

\begin{LeanCode}
L : A → B → C → D
\end{LeanCode}

当前目标是：

\begin{LeanCode}
⊢ D
\end{LeanCode}

那么：

\begin{LeanCode}
apply L
\end{LeanCode}

可能生成三个子目标：

\begin{LeanCode}
⊢ A
⊢ B
⊢ C
\end{LeanCode}

任务树：

\begin{Tree}
D
├── A
├── B
└── C
\end{Tree}

很多复杂数学定理本来就是：如果满足若干条件 \(A\)、\(B\)、\(C\)，就能推出结论 \(D\)。所以证明 \(D\) 时，可以先找有没有一个 lemma/theorem 能推出 \(D\)。如果有，就用 \texttt{apply} 把目标反推成 lemma 所需要的前提。

\section{第四类规则：改写和简化}

有些目标不需要拆成多个子目标，而是需要先把目标变简单。

\subsection{等式改写 rw}

如果已知：

\begin{LeanCode}
h : x = y
\end{LeanCode}

而目标中出现了 \(x\)，可以把 \(x\) 改成 \(y\)。

Lean 常用写法：

\begin{LeanCode}
rw [h]
\end{LeanCode}

如果要反向改写：

\begin{LeanCode}
rw [← h]
\end{LeanCode}

例子：

\begin{LeanCode}
theorem demo
  (x y : Nat)
  (h : x = y) :
  x + 1 = y + 1 := by
  rw [h]
\end{LeanCode}

目标变化：

\begin{Tree}
x + 1 = y + 1
↓ rw [h]
y + 1 = y + 1
\end{Tree}

然后就容易证明了。

\subsection{定义展开 unfold}

如果目标里有一个定义被包起来，可以展开它。

例如：

\begin{LeanCode}
def double (n : Nat) := n + n
\end{LeanCode}

目标：

\begin{LeanCode}
⊢ double n = n + n
\end{LeanCode}

可以：

\begin{LeanCode}
unfold double
\end{LeanCode}

直觉：有些复杂只是因为定义没有展开。展开后目标可能变得很简单。

\subsection{自动简化 simp}

\texttt{simp} 是 Lean 里非常常用的自动简化工具。

适合处理：
\begin{itemize}
  \item \(x + 0 = x\)
  \item \(\mathrm{True} \land P\)
  \item \(P \land \mathrm{True}\)
  \item \texttt{List.length []}
\end{itemize}

例子：

\begin{LeanCode}
theorem demo (x : Nat) : x + 0 = x := by
  simp
\end{LeanCode}

在拆分过程中，\texttt{simp} 很适合作为叶子目标的第一尝试：如果目标已经很简单，先试 \texttt{simp}。如果 \texttt{simp} 能关掉目标，就不需要继续拆。

\section{第五类规则：归纳和分类讨论}

当目标涉及递归结构时，通常要用 \texttt{induction} 或 \texttt{cases}。

\subsection{自然数归纳}

如果目标和自然数 \(n : \texttt{Nat}\) 有关，常常可以对 \(n\) 归纳。

目标形式：

\begin{LeanCode}
n : Nat
⊢ P n
\end{LeanCode}

归纳拆成：
\begin{itemize}
  \item base case：证明 \(P\ 0\)
  \item step case：假设 \(P\ n\)，证明 \(P(n+1)\)
\end{itemize}

Lean 常用写法：

\begin{LeanCode}
induction n with
| zero =>
  ...
| succ n ih =>
  ...
\end{LeanCode}

任务树：

\begin{Tree}
P n
├── P 0
└── 假设 P n，证明 P (n+1)
\end{Tree}

\subsection{列表归纳}

如果目标和列表有关：

\begin{LeanCode}
xs : List Nat
⊢ P xs
\end{LeanCode}

列表有两种情况：
\begin{enumerate}
  \item 空列表 \texttt{[]}；
  \item 非空列表 \texttt{x :: xs}。
\end{enumerate}

所以归纳拆成：
\begin{itemize}
  \item case nil：证明 \(P\ []\)
  \item case cons：假设 \(P\ xs\)，证明 \(P\ (x :: xs)\)
\end{itemize}

任务树：

\begin{Tree}
P xs
├── P []
└── 假设 P xs，证明 P (x :: xs)
\end{Tree}

\subsection{一般分类讨论}

如果某个对象只有有限几种形式，可以分类讨论。

例如：

\begin{LeanCode}
b : Bool
\end{LeanCode}

布尔值只有两种：
\begin{itemize}
  \item \texttt{true}
  \item \texttt{false}
\end{itemize}

所以：

\begin{LeanCode}
cases b
\end{LeanCode}

会拆成两个情况：
\begin{itemize}
  \item case true
  \item case false
\end{itemize}

直觉：如果一个对象本身只有几种可能，就把每种可能分别证明一遍。

\section{第六类规则：函数、结构、集合相等}

有些相等目标不能直接证明，需要先转化成更细的目标。

\subsection{函数相等 funext}

目标：

\begin{LeanCode}
⊢ f = g
\end{LeanCode}

如果 \(f\) 和 \(g\) 是函数，那么证明函数相等，通常要证明：对任意 \(x\)，\(f\ x = g\ x\)。

Lean 常用写法：

\begin{LeanCode}
funext x
\end{LeanCode}

目标变化：

\begin{Tree}
⊢ f = g
↓ funext x
x : α
⊢ f x = g x
\end{Tree}

直觉：两个函数相等，就是它们对每个输入的输出都相等。

\subsection{结构或集合相等 ext}

目标：

\begin{LeanCode}
⊢ S = T
\end{LeanCode}

如果 \(S\) 和 \(T\) 是集合，通常可以转化成：对任意 \(x\)，\(x \in S \leftrightarrow x \in T\)。

Lean 常用写法：

\begin{LeanCode}
ext x
\end{LeanCode}

直觉：两个集合相等，就是它们包含完全相同的元素。

\section{第七类规则：直接关闭目标}

有些目标已经足够小，不需要继续拆。

\subsection{exact}

如果目标正好已经在上下文里：

\begin{LeanCode}
hA : A
⊢ A
\end{LeanCode}

直接：

\begin{LeanCode}
exact hA
\end{LeanCode}

\subsection{rfl}

如果目标是自反等式：

\begin{LeanCode}
⊢ x = x
\end{LeanCode}

或者经过计算后左右两边完全一样，可以用：

\begin{LeanCode}
rfl
\end{LeanCode}

\subsection{contradiction}

如果上下文中有矛盾：

\begin{LeanCode}
hA : A
hnA : ¬ A
\end{LeanCode}

可以：

\begin{LeanCode}
contradiction
\end{LeanCode}

直觉：如果已知条件本身矛盾，那么当前目标可以被关闭。

\section{复杂证明如何递归拆分}

复杂 Lean 证明不是一次拆完的，而是不断看当前目标的最外层结构，然后递归拆下去。

例如目标是：

\begin{LeanCode}
⊢ A → ((B ∧ C) ∧ (D ↔ E))
\end{LeanCode}

第一步，最外层是 \(\to\)：

\begin{LeanCode}
intro hA
\end{LeanCode}

目标变成：

\begin{LeanCode}
hA : A
⊢ (B ∧ C) ∧ (D ↔ E)
\end{LeanCode}

第二步，新目标最外层是 \(\land\)：

\begin{LeanCode}
constructor
\end{LeanCode}

拆成：
\begin{itemize}
  \item 目标 1：\(\vdash B \land C\)
  \item 目标 2：\(\vdash D \leftrightarrow E\)
\end{itemize}

第三步，继续拆每个子目标：

\begin{Tree}
B ∧ C
├── B
└── C

D ↔ E
├── D → E
└── E → D
\end{Tree}

完整证明树：

\begin{Tree}
A → ((B ∧ C) ∧ (D ↔ E))
└── intro hA 后证明 (B ∧ C) ∧ (D ↔ E)
    ├── B ∧ C
    │   ├── B
    │   └── C
    └── D ↔ E
        ├── D → E
        └── E → D
\end{Tree}

如果继续拆 \(D \to E\)：

\begin{LeanCode}
intro hD
\end{LeanCode}

变成：

\begin{LeanCode}
hD : D
⊢ E
\end{LeanCode}

如果继续拆 \(E \to D\)：

\begin{LeanCode}
intro hE
\end{LeanCode}

变成：

\begin{LeanCode}
hE : E
⊢ D
\end{LeanCode}

所以真正的叶子目标是：

\begin{LeanCode}
⊢ B
⊢ C
hD : D ⊢ E
hE : E ⊢ D
\end{LeanCode}

\section{哪些拆分适合并行，哪些不适合}

如果把 Lean proof 看成任务树，有些分支可以并行证明，有些不能。

\subsection{适合并行的拆分}

\paragraph{\(A \land B\)}

\begin{Tree}
A ∧ B
├── A
└── B
\end{Tree}

证明 \(A\) 和证明 \(B\) 可以分开做。

\paragraph{\(A \leftrightarrow B\)}

\begin{Tree}
A ↔ B
├── A → B
└── B → A
\end{Tree}

两个方向可以分开做。

\paragraph{多前提 apply}

如果：

\begin{LeanCode}
L : A → B → C → D
\end{LeanCode}

目标是：

\begin{LeanCode}
⊢ D
\end{LeanCode}

用：

\begin{LeanCode}
apply L
\end{LeanCode}

变成：

\begin{LeanCode}
⊢ A
⊢ B
⊢ C
\end{LeanCode}

这些前提通常可以分开证明。

\subsection{不适合直接并行的拆分}

\paragraph{\(A \to B\)}

\begin{Tree}
A → B
└── 假设 A 后证明 B
\end{Tree}

这是单分支推进。

\paragraph{\(\forall x, P\ x\)}

\begin{Tree}
∀ x, P x
└── 任取 x 后证明 P x
\end{Tree}

这也是单分支推进。

\paragraph{\(\exists x, P\ x\)}

\begin{Tree}
∃ x, P x
├── 找 witness w
└── 证明 P w
\end{Tree}

证明 \(P\ w\) 依赖于先找到 \(w\)。

\paragraph{\(A \lor B\)}

\begin{Tree}
A ∨ B
├── 尝试证明 A
└── 尝试证明 B
\end{Tree}

这里只有一边需要成功。它更像竞争式搜索，不是两个任务都必须完成。

\section{最应该记住的核心规则}

如果只记核心规则，先记下面这些：

\begin{enumerate}
  \item 目标 \(A \land B\)：拆成证明 \(A\) 和证明 \(B\)。
  \item 目标 \(A \leftrightarrow B\)：拆成证明 \(A \to B\) 和 \(B \to A\)。
  \item 目标 \(A \to B\)：先假设 \(A\)，再证明 \(B\)。
  \item 目标 \(\forall x, P\ x\)：任取 \(x\)，再证明 \(P\ x\)。
  \item 目标 \(\exists x, P\ x\)：找 witness \(w\)，再证明 \(P\ w\)。
  \item 目标 \(A \lor B\)：证明 \(A\) 或证明 \(B\)，成功一边即可。
  \item 已知 \(h : A \land B\)：拆出 \(hA : A\) 和 \(hB : B\)。
  \item 已知 \(h : \exists x, P\ x\)：拆出 witness \(w\) 和证明 \(hw : P\ w\)。
  \item 已知 \(h : A \lor B\)：分情况讨论 \(A\) 和 \(B\)，两种情况都要证明目标。
  \item 已知 \(h : A \to B\)，目标是 \(B\)：用 \texttt{apply h}，把目标反推成 \(A\)。
  \item 目标里有等式或定义：用 \texttt{rw} / \texttt{unfold} / \texttt{simp} 先化简。
  \item 目标涉及 \texttt{Nat}、\texttt{List}、递归结构：用 \texttt{induction} 或 \texttt{cases} 拆成基本情况和递归情况。
  \item 目标已经很小：用 \texttt{exact} / \texttt{rfl} / \texttt{simp} / \texttt{contradiction} 关闭。
\end{enumerate}

\section{一句话总结}

Lean 证明拆分的核心不是死记 tactic，而是看结构：

\begin{quote}
目标是什么形状，就用对应规则把它变成更小的目标；已知条件是什么形状，就把它拆成更可用的条件；这个过程递归进行，直到所有叶子目标都能被 \texttt{exact}、\texttt{simp}、\texttt{rfl}、\texttt{apply}、\texttt{rw}、\texttt{induction} 等方式解决。
\end{quote}

复杂 Lean 证明可以理解为：

\begin{Tree}
逻辑结构递归拆分
+
叶子目标逐个关闭
=
完整 Lean proof tree
\end{Tree}

\end{document}
